\section{Method}
\label{sec:method}

% TODO: write the methodology. Pipeline is Detect -> Embed -> Type.
% Skeleton updated to the current approach (V2, 2026-06-11). The old GMM-per-label /
% anchor-word / OOD / hierarchy design is NOT the method anymore; the GMM survives
% only as one ablation baseline of the typing head.
%
% \subsection{Overview}
%   - Pipeline figure: Text -> MentionDetector -> Matryoshka Embedder -> Typing head -> label.
%   - Three design principles:
%     (i)   no heavy encoder is trained from scratch (off-the-shelf components);
%     (ii)  a light contrastive fine-tuning, refined by error-driven hard-negative
%           mining, makes a general-purpose embedding discriminative for entity types;
%     (iii) Matryoshka-truncatable embeddings give a compute/quality knob.
%
% \subsection{Mention Detection}
%   - GLiNER (frozen, off-the-shelf): raw text + candidate labels -> spans.
%   - Justify: strong zero-shot detector, no retraining; it is the shared end-to-end bottleneck.
%
% \subsection{Matryoshka Entity Embedding}
%   - Nomic v1.5, span-in-context, truncate_dim in {768, 512, 256, 128, 64}.
%   - e_i = Embed(x_i; D) in R^D.
%
% \subsection{Contrastive Fine-tuning with Hard-Negative Mining}
%   - Triplet margin loss (margin 1) on CoNLL-2003 source spans.
%   - Second pass: fine-tune again on the entity pairs the model most often confuses
%     (hard negatives in the sense of FaceNet, \cite{schroff2015facenet}), mixed with
%     correctly-typed mentions for diversity; errors are mined per training split by
%     cross-validation, so no test label is seen.
%
% \subsection{Entity Typing}
%   - Supervised: a balanced linear classifier (LinearSVC, class_weight=balanced) fitted
%     on each target's training labels; give the n_total / (n_classes * count) weighting.
%   - Zero-shot variant (OPENER-ZS): no target labels; assign each mention to the nearest
%     label-name prototype (centroid of the label's anchor words embedded in the
%     contrastive space) by cosine similarity.
%   - A per-class GMM is reported only as an ablation baseline (\autoref{tab:ablation}).
%
% \subsection{Evaluation protocol}
%   - Offset alignment + FN/FP sentinels; AMI; typing-on-gold vs end-to-end; latency + energy.

\emph{[Placeholder - method to be written once the configuration is frozen.]}
